% Part: computability
% Chapter: computability-theory
% Section: equiv-ce-defs

\documentclass[../../../include/open-logic-section]{subfiles}

\begin{document}

\olfileid[hi]{cmp}{thy}{eqc}

\olsection[C. E. समुच्चयों की परिभाषाएँ]{संगणनीयतः परिगणनीय समुच्चयों
  की समतुल्य परिभाषाएँ}


निम्न प्रमेय संगणनीयतः परिगणनीय होने के अर्थ के कई महत्त्वपूर्ण समतुल्य
कथन देता है।

\begin{thm}
\ollabel{thm:ce-equiv}
$S$ को प्राकृतिक संख्याओं का समुच्चय मानें। तब निम्नलिखित कथन समतुल्य हैं:
\begin{enumerate}
\item\ollabel{case:ce} $S$ संगणनीयतः परिगणनीय है।
\item\ollabel{case:ran-pc} $S$ किसी \emph{आंशिक} संगणनीय फलन का परिसर है।
\item\ollabel{case:ran-prim} $S$ रिक्त है अथवा किसी आदिम पुनरावर्ती फलन का परिसर है।
\item\ollabel{case:ce-domain} $S$ किसी आंशिक संगणनीय फलन का \emph{प्रांत} है।
\end{enumerate}
\end{thm}

\begin{explain}
पहले तीन उपवाक्य कहते हैं कि किसी भी अरिक्त संगणनीयतः परिगणनीय समुच्चय
को समतुल्य रूप से किसी संगणनीय फलन, किसी आंशिक संगणनीय फलन अथवा किसी
आदिम पुनरावर्ती फलन द्वारा परिगणित माना जा सकता है। चौथा उपवाक्य बताता
है कि यदि $S$ संगणनीयतः परिगणनीय है, तो किसी सूचकांक~$e$ के लिए
\[
S = \Setabs{x}{\cfind{e}(x) \fdefined}.
\]
दूसरे शब्दों में, $S$ उन आगतों का समुच्चय है जिन पर $\cfind{e}$ की
संगणना रुकती है। इसी कारण संगणनीयतः परिगणनीय समुच्चयों को कभी-कभी
\emph{अर्ध-निर्णेय} कहते हैं: यदि कोई संख्या समुच्चय में है, तो अंततः
``हाँ'' मिल जाती है; लेकिन यदि वह उसमें नहीं है, तो कभी ``नहीं'' नहीं मिलती!{}
\end{explain}

\begin{proof}
चूँकि प्रत्येक आदिम पुनरावर्ती फलन संगणनीय और प्रत्येक संगणनीय फलन
आंशिक संगणनीय है, \olref{case:ran-prim} से \olref{case:ce} तथा
\olref{case:ce} से~\olref{case:ran-pc} निष्पन्न होता है। (ध्यान दें कि
यदि $S$ रिक्त है, तो $S$~उस आंशिक संगणनीय फलन का परिसर है जो कहीं भी
परिभाषित नहीं है।) यदि हम दिखा दें कि \olref{case:ran-pc} से
\olref{case:ran-prim} निष्पन्न होता है, तो पहले तीनों उपवाक्यों की
समतुल्यता सिद्ध हो जाएगी।

अतः मान लें कि $S$, आंशिक संगणनीय फलन $\cfind{e}$ का परिसर है। यदि $S$
रिक्त है, तो कार्य पूरा हुआ। अन्यथा $a$ को~$S$ का कोई अवयव मानें। क्लीन
के प्रसामान्य रूप प्रमेय से हम लिख सकते हैं
\[
\cfind{e}(x) = U(\umin{s}{T(e, x, s)}).
\]
विशेषतः, $\cfind{e}(x) \fdefined$ और $= y$ तभी और केवल तभी जब कोई $s$
ऐसा हो कि $T(e, x, s)$ तथा $U(s) = y$। $f(z)$ को परिभाषित करें:
\[
f(z) = \begin{cases}
  U((z)_1) & \text{यदि $T(e, (z)_0, (z)_1)$} \\
  a        & \text{अन्यथा।}
\end{cases}
\]
तब $f$ आदिम पुनरावर्ती है, क्योंकि $T$ और $U$ ऐसे ही हैं।
\iftag{TMs}{ट्यूरिंग मशीनों के पदों में, यदि $z$ किसी ऐसे युग्म
  $\tuple{(z)_0, (z)_1}$ को कूटित करता है जिसमें $(z)_1$, मशीन~$M_e$ की
  आगत $(z)_0$ पर रुकने वाली संगणना है, तो $f$ संगणना का निर्गत लौटाता
  है; अन्यथा वह~$a$ लौटाता है।}

हमें दिखाना है कि $S$,~$f$ का परिसर है; अर्थात् किसी भी प्राकृतिक
संख्या~$y$ के लिए $y \in S$ तभी और केवल तभी जब वह~$f$ के परिसर में है।
आगे की दिशा में मान लें कि $y \in S$। तब $y$, $\cfind{e}$ के परिसर में
है, अतः किन्हीं $x$ और~$s$ के लिए $T(e,x,s)$ सत्य है तथा $U(s) = y$।
पर तब $y = f(\tuple{x,s})$। उलटी दिशा में मान लें कि $y$,~$f$ के परिसर
में है। तब या तो $y = a$, अथवा किसी~$z$ के लिए $T(e,(z)_0,(z)_1)$ और
$U((z)_1) = y$। चूँकि बाद वाली स्थिति में $\cfind{e}(x) \fdefined =
y$,
दोनों ही स्थितियों में $y$,~$S$ में है।

(संकेत $\cfind{e}(x) \fdefined = y$ का अर्थ है कि ``$\cfind{e}(x)$
परिभाषित और $y$ के बराबर है।'' हम $\cfind{e}(x) =
y$ भी लिख सकते थे,
पर अतिरिक्त तीर कभी-कभी यह याद दिलाने में सहायक है कि हम किसी आंशिक
फलन से व्यवहार कर रहे हैं।)

\olref{thm:ce-equiv} का प्रमाण पूरा करने के लिए यह दिखाना पर्याप्त है कि
\olref{case:ce} और~\olref{case:ce-domain} समतुल्य हैं। पहले दिखाएँ कि
\olref{case:ce} से~\olref{case:ce-domain} निष्पन्न होता है। मान लें कि
$S$ किसी संगणनीय फलन~$f$ का परिसर है, अर्थात्
\[
S = \Setabs{y}{\text{किसी $x$ के लिए, } f(x) = y}.
\]
मानें
\[
g(y) = \umin{x}{(f(x) = y)}.
\]
तब $g$ एक आंशिक संगणनीय फलन है, और $g(y)$ तभी और केवल तभी परिभाषित है
जब किसी~$x$ के लिए $f(x) = y$। दूसरे शब्दों में,~$g$ का प्रांत~$f$ का
परिसर है। \iftag{TMs}{ट्यूरिंग मशीनों के पदों में: ट्यूरिंग मशीन~$F$
मिलने पर, जो~$S$ के अवयवों का परिगणन करती है, $G$ वह ट्यूरिंग मशीन हो जो
$S$ को~$F$ के निर्गतों में किसी दिए अवयव को खोजकर अर्ध-निर्णीत करती है;
अवयव मिलने पर रुकती है और न मिलने पर सदा खोजती रहती है।}{}

अंततः, \olref{case:ce-domain} से~\olref{case:ce} दिखाने के लिए मान लें
कि $S$~आंशिक संगणनीय फलन~$\cfind{e}$ का प्रांत है, अर्थात्
\[
S = \Setabs{x}{\cfind{e}(x) \fdefined}.
\]
यदि $S$ रिक्त है, तो कार्य पूरा हुआ; अन्यथा $a$ को~$S$ का कोई अवयव मानें।
$f$ को परिभाषित करें:
\[
f(z) = \begin{cases}
(z)_0 & \text{यदि $T(e,(z)_0,(z)_1)$} \\
a & \text{अन्यथा।}
\end{cases}
\]
तब, ऊपर की तरह, कोई संख्या $x$,~$f$ के परिसर में तभी और केवल तभी है जब
$\cfind{e}(x) \fdefined$, अर्थात् तभी और केवल तभी जब $x \in S$।
\iftag{TMs}{ट्यूरिंग मशीनों के पदों में: मशीन $M_e$ मिलने पर, जो
$S$ को अर्ध-निर्णीत करती है, सभी संभव ट्यूरिंग मशीन संगणनाओं को चलाकर और रुकने
वाली संगणनाओं के अनुरूप आगत लौटाकर $S$ के अवयवों का परिगणन करें।}{}
\end{proof}

\olref{thm:ce-equiv} का उपवाक्य~\olref{case:ce-domain} संगणनीयतः
परिगणनीय समुच्चयों के परिगणन का सुविधाजनक ढंग देता है: प्रत्येक~$e$ के
लिए $W_e$, $\cfind{e}$ के प्रांत को निरूपित करे, अर्थात्
\[
W_e = \Setabs{x}{\cfind{e}(x) \fdefined}.
\] 
तब यदि $A$ कोई संगणनीयतः परिगणनीय समुच्चय है, तो $A = W_e$, किसी~$e$ के लिए।

अगला प्रमेय संगणनीयतः परिगणनीय समुच्चयों का एक और अभिलक्षण देता है।

\begin{thm}
\ollabel{thm:exists-char}
कोई समुच्चय $S$ तभी और केवल तभी संगणनीयतः परिगणनीय है जब कोई संगणनीय
संबंध $R(x,y)$ ऐसा हो कि
\[
S = \Setabs{ x }{ \lexists[y][R(x,y)] }.
\]
\end{thm}

\begin{proof}
आगे की दिशा में मान लें कि $S$ संगणनीयतः परिगणनीय है। तब किसी $e$ के
लिए $S = W_e$।~$e$ के इस मान पर हम $S$ को इस प्रकार लिख सकते हैं:
\[
S = \Setabs{ x }{ \lexists[y][T(e, x, y)] }.
\]
उलटी दिशा में मान लें कि $S = \Setabs{ x }{
  \lexists[y][R(x, y)] }$। $f$~को परिभाषित करें:
\[
f(x) \simeq \umin{y}{R(x, y)}.
\]
तब $f$ आंशिक संगणनीय है और $S$,~$f$ का प्रांत है।
\end{proof}


\end{document}
